\documentclass[11pt,a4paper]{article}
\usepackage{fontspec}
\setmainfont{LibertinusSerif-Regular.otf}[
 BoldFont=LibertinusSerif-Bold.otf,
 ItalicFont=LibertinusSerif-Italic.otf,
 BoldItalicFont=LibertinusSerif-BoldItalic.otf]
\setmonofont{lmmono10-regular.otf}
\usepackage[russian]{babel}
\usepackage[margin=25mm]{geometry}
\usepackage{amsmath,amssymb,amsthm}
\usepackage{booktabs,array,longtable}
\usepackage{xurl}
\usepackage[unicode,hidelinks]{hyperref}
\usepackage{microtype}
\setlength{\emergencystretch}{2em}
\setlength{\parskip}{3pt}
\setlength{\parindent}{1.3em}
\widowpenalty=10000
\renewcommand{\arraystretch}{1.2}
\numberwithin{equation}{section}
\newtheorem{theorem}{Теорема}[section]
\newtheorem{lemma}[theorem]{Лемма}
\newtheorem{corollary}[theorem]{Следствие}
\theoremstyle{definition}
\newtheorem{definition}{Определение}
\theoremstyle{remark}
\newtheorem{remark}[theorem]{Замечание}
\newcommand{\QQ}{\mathbb{Q}}
\newcommand{\HH}{\mathbb{H}_{\QQ}}
\newcommand{\OO}{\mathbb{O}_{\QQ}}
\newcommand{\J}{\mathcal{J}}
\newcommand{\D}{\mathcal{D}}
\newcommand{\Ham}{\mathrm{d}_{\mathrm H}}
\newcommand{\diag}{\operatorname{diag}}
\newcommand{\wt}{\operatorname{wt}}
\newcommand{\leanref}[2]{\par\smallskip\noindent{\small\textit{Проверено в Lean:} \nolinkurl{#1}; \nolinkurl{#2}.}\par}
\hypersetup{pdftitle={Связи математических объектов и общие структуры: три пары с проверкой в Lean},pdfauthor={}}
\title{Связи математических объектов\newline и общие структуры:\newline
три пары с проверкой в Lean}
\author{\small Препринт, версия 1.5}
\date{29 октября 2026 года}

\begin{document}
\maketitle
\begin{abstract}
Основная задача работы --- выбрать связанные математические объекты,
установить точную связь между ними и выяснить, можно ли выделить
общую структуру более высокого уровня. Для этого рассматриваются
три пары объектов. Сначала каждый объект определяется отдельно,
затем задаётся и доказывается соответствие между объектами пары.
На основе доказанной связи выделяются общие данные и оценивается
кандидат на <<суперобъект>>. Шесть определений охватывают операторную
и геометрическую модели, модель умножения и спектральную модель,
отношение нулевых ассоциаторов и геометрический код. В этих примерах
общими структурами выступают веса расстояний, взвешенное разложение
пространства и инцидентность с отображениями наблюдения.
Содержательность связей подтверждается теоремами о восстановлении,
энергии остатка и устойчивости к ошибкам, проверенными в Lean~4.19.0.
Рассмотренные пары иллюстрируют общую постановку для математических
объектов. Историческая новизна и высокая внутренняя сложность
выделенных структур отдельно не установлены.
\end{abstract}
\noindent\textbf{Ключевые слова:} октонионы, алгебра Альберта, ассоциатор,
геометрия расстояний, спектральная оценка, конечные коды, Lean.

\section{Общая задача: от связанных объектов к общей структуре}
\subsection{Цель исследования}
Требуется выбрать математические объекты, найти и доказать связь
между ними, а затем выяснить, какие общие данные объясняют эту связь.
Если такие данные образуют самостоятельную структуру с согласованными
представлениями в исходных объектах, её можно рассматривать как
кандидата на объект более высокого уровня, или \emph{суперобъект}.

Эта постановка относится к математическим объектам разных видов:
пространствам, операторам, геометриям, отношениям, кодам и другим
структурам. Конкретные пары в работе выбраны для проверки этой идеи
на примерах, для которых можно явно задать оба объекта и доказать
соответствие между ними. Они не ограничивают общую задачу одним
разделом математики.

Для каждой пары используется одна последовательность:
\begin{enumerate}
 \item Определить первый математический объект.
 \item Независимо указать, что представляет собой второй объект.
 \item Задать связь между ними и установить, что она сохраняет.
 \item Доказать заявленное соответствие.
 \item Выделить общую структуру и сформулировать, в каком смысле
 она может рассматриваться как кандидат на суперобъект.
\end{enumerate}
Таким образом, исходной задачей является поиск и проверка связей
между объектами. Восстановление, спектральные оценки и исправление
ошибок ниже служат конкретными способами проверить найденные связи.

\subsection{Три пары и шесть определений}
Сначала рассматриваются следующие пары; точные определения даны
в соответствующих разделах.
\begin{center}\small
\begin{tabular}{@{}p{.07\textwidth}p{.42\textwidth}p{.42\textwidth}@{}}
\toprule
Пара & Первый объект & Второй объект\\
\midrule
1 & Семейство операторов ассоциатора
(определение~\ref{def:operators}) & Геометрия размеченных троек точек
(определение~\ref{def:points})\\
2 & Пространство с оператором умножения и квадратичной формой
(определение~\ref{def:multiplication}) & Трёхблочная спектральная модель
(определение~\ref{def:spectral})\\
3 & Отношение нулевых ассоциаторов
(определение~\ref{def:incidence}) & Код размеченных плоскостей Фано
(определение~\ref{def:fano-code})\\
\bottomrule
\end{tabular}
\end{center}
В первой паре выясняется, какие геометрические данные определяют
оператор и восстанавливаются по нему. Во второй устанавливается
соответствие между действием умножения и взвешенными спектральными
блоками. В третьей проверяется, что отношение на тройках и кодовое
слово представляют одну инцидентность, которую можно восстановить
по части наблюдений.

\subsection{Когда можно выделить объект более высокого уровня}
Наличие связи само по себе ещё не определяет суперобъект.
Нужно явно указать общие данные, способы получения из них обоих
представлений и свойства, согласуемые этими способами. В дальнейшем
именно такой набор данных и представлений называется кандидатом
на суперобъект. Вывод для каждой пары делается после установления
связи и относится к конкретно сохранённой структуре.

Термин <<более высокого уровня>> здесь описывает роль общей структуры
по отношению к её представлениям. Высокая внутренняя сложность,
универсальность и историческая новизна являются отдельными
исследовательскими вопросами; доказательство соответствия само по
себе их не устанавливает.

\subsection{Место результатов в литературе}
Связанные идеи классические: коэффициенты Пирса и операторы Шура
рассматриваются в~\cite{gowda}, восстановление по размеченным расстояниям
--- в~\cite{dokmanic}, спектр осей Monster-типа --- в~\cite{mcinroy}.
Октонная интерпретация плоскости Фано описана в~\cite{baez}, а семейство
из 30 размеченных плоскостей и кодовые конструкции на 35 позициях
--- в~\cite{liebler}. Конкретные сокращённые наборы тестов, приведённые
ниже, не были отождествлены с результатом в просмотренных источниках;
это не является доказательством их исторического приоритета.
Повторное сопоставление с первичными источниками 16 октября 2026 года
показало, что первые две композиции являются специализациями известных
формул и оценок; граница $N_{\min}\ge9$ также следует из опубликованной
общей теоремы. Поэтому работа представляет явное согласование моделей
и формальные сертификаты, а не установленное открытие ранее неизвестных
связей. Точные основания этого вывода приведены после соответствующих
доказательств.

\section{Обозначения для выбранных пар}\label{sec:model}
Все скаляры принадлежат $\QQ$. Пусть $\HH$ --- рациональные кватернионы
с обычным сопряжением. Определим октонные координаты
$\OO=\HH\times\HH$ с конвенцией Кэли--Диксона
\begin{equation}\label{eq:oct}
 (a,b)(c,d)=(ac-d\bar b,\;\bar a d+cb),
 \qquad \overline{(a,b)}=(\bar a,-b).
\end{equation}
Скаляр $r$ отождествляется с $((r,0,0,0),0)$.
Порядок множителей в~\eqref{eq:oct} является частью определения;
все дальнейшие вычисления используют именно эту конвенцию.

Элемент $x\in\J$ задаётся 27 рациональными координатами:
\begin{equation}\label{eq:matrix}
 x=\begin{pmatrix}
 x_0 & u_{01} & u_{02}\\
 \bar u_{01} & x_1 & u_{12}\\
 \bar u_{02} & \bar u_{12} & x_2
 \end{pmatrix},
 \qquad x_i\in\QQ,\quad u_{ij}\in\OO.
\end{equation}
На этих координатах задано симметризованное произведение
\begin{equation}\label{eq:jordan}
 x\circ y=\tfrac12(xy+yx),\qquad
 (xy)_{ij}=\sum_{k=0}^{2}x_{ik}y_{kj}.
\end{equation}
Каждое слагаемое содержит только одно октонное умножение, поэтому
бинарный матричный продукт определён однозначно. В произведениях
трёх элементов скобки сохраняются явно.

В Lean координатное вложение в матрицы инъективно; доказаны эрмитовость
симметризованного произведения и его совместимость с координатной
операцией. Это конкретная рациональная модель стандартных координат
Альберта. Полное тождество Иордана и классификационные свойства
алгебры Альберта в рассматриваемой формализации отдельно не доказаны.
Доказательства ниже используют явную операцию~\eqref{eq:jordan}.

Для $u\in\OO$ обозначим через $|u|_2^2$ сумму квадратов его восьми
рациональных координат. Введём квадратичную форму
\begin{equation}\label{eq:size}
 S(x)=x_0^2+x_1^2+x_2^2+
 2\bigl(|u_{01}|_2^2+|u_{02}|_2^2+|u_{12}|_2^2\bigr).
\end{equation}
В дальнейших оценках используется эта явная сумма квадратов;
не требуется постулировать абстрактную норму.

\section{Первая пара: операторы и геометрия точек}\label{sec:geometry}
\subsection{Определения двух объектов}
\begin{definition}[Семейство операторов ассоциатора]\label{def:operators}
Первый объект --- семейство $\mathcal A_1=\{\D_d:d\in\QQ^3\}$
операторов на $\J$. Для $d=(d_0,d_1,d_2)$ и $a=\diag(d)$ положим
\begin{equation}\label{eq:defect}
 \D_d(x)=(a\circ a)\circ x-a\circ(a\circ x).
\end{equation}
Два элемента семейства считаются равными, если совпадают их действия
на всех $x\in\J$.
\end{definition}

\begin{definition}[Геометрия размеченных троек]\label{def:points}
Второй объект --- множество $\mathcal B_1$ конфигураций трёх
размеченных точек рациональной прямой с точностью до движения прямой.
Конфигурация задаётся $d=(d_0,d_1,d_2)\in\QQ^3$; тройки $d,e$
представляют одну конфигурацию, если
$e_i=\varepsilon d_i+c$ для всех $i$, где
$\varepsilon\in\{1,-1\}$ и $c\in\QQ$.
Метки сохраняются, совпадения точек допускаются.
\end{definition}

\subsection{Какую связь устанавливаем}
Каждой тройке $d$ сопоставлен оператор $\D_d$. Требуется проверить,
что это соответствие устанавливает взаимно однозначную связь
между конфигурациями $\mathcal B_1$ и операторами $\mathcal A_1$.
Для этого нужно доказать, что один и тот же оператор получается
ровно у троек, отличающихся общим сдвигом или отражением.

Сопоставление проводится через три числа
\[
 w_{ij}(d)=\tfrac14(d_i-d_j)^2,\qquad 0\le i<j\le2.
\]
Для второго объекта это четверти квадратов расстояний. Для первого
они должны оказаться коэффициентами действия оператора.
Лемма~\ref{lem:weights} и теорема~\ref{thm:rigidity} устанавливают
это соответствие; затем графовая интерпретация даёт ещё одно
представление тех же данных.

\subsection{Утверждение и доказательство}

\begin{lemma}\label{lem:weights}
Оператор $\D_d$ равен нулю на диагональных координатах и умножает
каждую из восьми координат блока $u_{ij}$ на $w_{ij}(d)$.
\end{lemma}
\begin{proof}
Непосредственно из~\eqref{eq:jordan} следует
\[
 (a\circ x)_i=d_i x_i,\qquad
 (a\circ x)_{ij}=\tfrac12(d_i+d_j)u_{ij}.
\]
При этом $a\circ a=\diag(d_0^2,d_1^2,d_2^2)$.
На диагонали два слагаемых в~\eqref{eq:defect} совпадают.
На блоке $ij$ их разность имеет коэффициент
\[
 \frac{d_i^2+d_j^2}{2}-\frac{(d_i+d_j)^2}{4}
 =\frac{(d_i-d_j)^2}{4}.
\]
Вычисление действует на произвольный октонный блок, а не только
на скалярные элементы.
\end{proof}

\begin{theorem}[Восстановление размеченных параметров]\label{thm:rigidity}
Для любых $d,e\in\QQ^3$ выполняется
\begin{equation}\label{eq:rigidity}
 \D_d=\D_e\quad\Longleftrightarrow\quad
 \bigl(\exists c\in\QQ\ \forall i:\ e_i=d_i+c\bigr)
 \ \lor\
 \bigl(\exists c\in\QQ\ \forall i:\ e_i=-d_i+c\bigr).
\end{equation}
Равенство слева есть равенство операторов на всех координатах $\J$.
Совпадающие значения $d_i$ допускаются.
\end{theorem}
\begin{proof}
Подставляя в лемму~\ref{lem:weights} по одному скалярному элементу
в каждый внедиагональный блок, получаем: равенство операторов
эквивалентно равенству трёх размеченных квадратов расстояний.

Положим $p=d_1-d_0$, $q=d_2-d_0$,
$p'=e_1-e_0$, $q'=e_2-e_0$. Тогда
\[
 (p')^2=p^2,\qquad(q')^2=q^2,\qquad(p'-q')^2=(p-q)^2.
\]
Над $\QQ$ первое и второе равенства дают
$p'=\varepsilon p$, $q'=\delta q$, где
$\varepsilon,\delta\in\{1,-1\}$. При одинаковых знаках получаем
общий сдвиг либо отражение с общим сдвигом. При различных знаках
третье равенство даёт $4pq=0$. Значит, $p=0$ или $q=0$.
Знак соответствующей нулевой координаты можно заменить, согласовав
его с другим знаком. Таким образом, и в вырожденном случае существует
единый $\varepsilon\in\{1,-1\}$, для которого
$e_i-e_0=\varepsilon(d_i-d_0)$ при всех $i$.
Достаточно взять $c=e_0-\varepsilon d_0$.

Обратно, общий сдвиг и отражение сохраняют квадраты всех разностей,
а значит, по лемме~\ref{lem:weights}, весь оператор.
\end{proof}
\leanref{AlbertSpectralGeometry.lean}{diagonalDefect_eq_iff_translation_or_reflection}

\subsection{Графовая интерпретация}
Рассмотрим полный граф на трёх размеченных вершинах с весами $w_{ij}(d)$.
Его взвешенный лапласиан $L_d$ имеет диагональные элементы
$\sum_{j\ne i}w_{ij}$ и внедиагональные элементы $-w_{ij}$.
Для любого $f=(f_0,f_1,f_2)\in\QQ^3$
\begin{equation}\label{eq:energy}
 \mathcal E_d(f)=\sum_{i<j}w_{ij}(d)(f_i-f_j)^2
 =f^{\mathsf T}L_df\ge0.
\end{equation}
Равенство получается раскрытием квадратов, неотрицательность --- из
$w_{ij}\ge0$. Если вложить скалярный градиент $f_i-f_j$ в октонные
рёбра, та же сумма выражает его спаривание с действием $\D_d$.
Каждое ребро в~\eqref{eq:energy} учитывается один раз.
Операторы $\D_d$ и $L_d$ действуют на разных пространствах:
первый --- на координатах~\eqref{eq:matrix}, второй --- на вершинах.
Их равенство не утверждается. Совпадение только спектров $\D_d$
также не заменяет гипотезу теоремы~\ref{thm:rigidity}.
\leanref{AlbertSpectralGeometry.lean}{associatorDirichlet_graphLaplacian}

\paragraph{Сопоставление с литературой.}
Формула действия $a\circ x$ на блоке $ij$ является стандартной
формулой Пирса~\cite[формула (2), с.~720]{gowda2012}.
Вычитание квадратов её коэффициентов даёт лемму~\ref{lem:weights}.
Восстановление по всем размеченным расстояниям с точностью до движения
есть общий принцип геометрии расстояний~\cite[формулы (7), (10)]{dokmanic};
теорема~\ref{thm:rigidity} даёт его элементарную рациональную одномерную
специализацию, включая совпадения точек. Графовое
тождество~\eqref{eq:energy} известно в общем виде, например
из~\cite[с.~920]{kalofolias}. Таким образом, составная конструкция
не даёт оснований заявлять новую математическую связь.

\subsection{Конкретный пример и смысл связи}
Возьмём $d=(0,2,6)$. Тогда $w(d)=(1,9,4)$ и
\[
 L_d=\begin{pmatrix}10&-1&-9\\-1&5&-4\\-9&-4&13\end{pmatrix}.
\]
Оператор $\D_d$ умножает октонные блоки $u_{01},u_{02},u_{12}$
соответственно на $1,9,4$. Для $f=(1,0,-1)$ получаем
\[
 \mathcal E_d(f)=1\cdot1^2+9\cdot2^2+4\cdot1^2
 =41=f^{\mathsf T}L_df.
\]
Тройки $(5,7,11)=d+5$ и $(8,6,2)=-d+8$ дают те же веса,
оператор и лапласиан. Теорема~\ref{thm:rigidity} доказывает,
что любой другой набор с тем же оператором получается именно
сдвигом или отражением со сдвигом.

\subsection{Какая общая структура выделена}
После доказанного соответствия общими данными пары становятся
размеченные веса $w(d)=(w_{01},w_{02},w_{12})$.
По ним определяется действие $\D_d$ и восстанавливается конфигурация
точек. Те же веса задают графовую энергию. Поэтому набор весов
вместе с этими представлениями является явным кандидатом на
суперобъект для первой пары.

Здесь выделена общая структура конкретного семейства операторов
и конфигураций. Соответствие не отождествляет всю модель $\J$
с графом: связываются именно данные, описанные в
определениях~\ref{def:operators} и~\ref{def:points}.

\section{Вторая пара: умножение и спектральная модель}\label{sec:spectral}
\subsection{Определения двух объектов}
\begin{definition}[Пространство с оператором умножения]\label{def:multiplication}
Первый объект --- $\mathcal A_2=(\J,M_P,S)$, где
$P=\diag(1,0,0)$, $M_P(x)=P\circ x$, а $S$ задана
в~\eqref{eq:size}. В этой паре рассматриваются выделенный оператор
и квадратичная форма на пространстве $\J$.
\end{definition}

\begin{definition}[Трёхблочная спектральная модель]\label{def:spectral}
Второй объект --- $\mathcal B_2=(V,T_B,S_B)$, где
\[
 V=\QQ\oplus\OO^2\oplus(\QQ^2\oplus\OO).
\]
Для $z=(\alpha,(u,v),(\beta,\gamma,w))\in V$ зададим
\begin{align*}
 T_Bz&=(\alpha,(u/2,v/2),(0,0,0)),\\
 S_B(z)&=\alpha^2+2(|u|_2^2+|v|_2^2)
       +\beta^2+\gamma^2+2|w|_2^2.
\end{align*}
Таким образом, $T_B$ действует на трёх блоках коэффициентами
$1$, $1/2$, $0$, а $S_B$ измеряет сумму их квадратичных вкладов.
\end{definition}

\subsection{Какую связь устанавливаем}
Сопоставим элементу $x\in\J$ из~\eqref{eq:matrix} его координаты,
сгруппированные следующим образом:
\[
 U(x)=\bigl(x_0,(u_{01},u_{02}),(x_1,x_2,u_{12})\bigr)\in V.
\]
Это обратимое линейное перегруппирование координат. Формула умножения на $P$
и определение $S$ дают
\[
 U(M_Px)=T_B(Ux),\qquad S_B(Ux)=S(x).
\]
Действительно, $P\circ x$ сохраняет $x_0$, делит $u_{01},u_{02}$
на два и обращает остальные координаты в нуль; выражение $S_B(Ux)$
совпадает с~\eqref{eq:size}. Значит, соответствие сохраняет выделенный
оператор и квадратичную форму. Сохранение всего произведения
$x\circ y$ этой парой утверждений не задаётся.

Чтобы установить следствие этой связи для произвольного элемента,
рассмотрим остаток $R_\lambda x=M_Px-\lambda x$. Тогда
\[
 S(R_\lambda x)=S_B\bigl((T_B-\lambda I)Ux\bigr).
\]
Энергия остатка определяется тремя спектральными блоками размерностей
$1$, $16$ и $10$. Ниже эта формула раскрывается в точное тождество
и острую нижнюю оценку. Связь с первой парой выражается также
равенством $\D_{(1,0,0)}=M_P-M_P^2$, поскольку $P\circ P=P$.

\subsection{Утверждение и доказательство}

\begin{theorem}[Точная формула и острая оценка]\label{thm:residual}
Для всех $x\in\J$ и $\lambda\in\QQ$ имеем
\begin{equation}\label{eq:residual}
 S(R_\lambda x)-\lambda^2S(x)
 =(1-2\lambda)x_0^2+
 2\left(\tfrac14-\lambda\right)
 \bigl(|u_{01}|_2^2+|u_{02}|_2^2\bigr).
\end{equation}
При $\lambda\le1/4$ отсюда следует
\begin{equation}\label{eq:bound}
 S(R_\lambda x)\ge\lambda^2S(x).
\end{equation}
Для каждого такого $\lambda$ константу $\lambda^2$ нельзя увеличить
в оценке, требуемой для всех $x\in\J$.
\end{theorem}
\begin{proof}
Положим
\[
 a=x_0^2,\qquad b=|u_{01}|_2^2+|u_{02}|_2^2,
 \qquad c=x_1^2+x_2^2+2|u_{12}|_2^2.
\]
Тогда $S(x)=a+2b+c$, а формула действия $P$ даёт
\[
 S(R_\lambda x)=(1-\lambda)^2a+
 2\left(\tfrac12-\lambda\right)^2b+\lambda^2c.
\]
Вычитание $\lambda^2(a+2b+c)$ доказывает~\eqref{eq:residual}.
При $\lambda\le1/4$ оба оставшихся коэффициента неотрицательны.
Для $x_* =\diag(0,1,0)$ имеем $S(x_*)=1$ и
$S(R_\lambda x_*)=\lambda^2$, что доказывает остроту.
\end{proof}
\leanref{SpectralDefect.lean}{residual_sizeSq; residual_lower_bound; sharpWitness_residual_sizeSq}

\paragraph{Сопоставление с литературой.}
Общая оценка расстояния до спектра самосопряжённого оператора
имеет вид
\[
 \|(T-\lambda I)x\|\ge
 \operatorname{dist}(\lambda,\operatorname{spec}T)\,\|x\|;
\]
см.~\cite[Theorem~6.1, с.~51]{pankrashkin}.
После расширения скаляров до $\mathbb R$ наш $M_P$ самосопряжён
для скалярного произведения с квадратом нормы $S$.
При $\lambda\le1/4$ ближайшая точка его спектра
$\{0,1/2,1\}$ --- $0$ (при $1/4$ имеется равенство расстояний),
поэтому общий результат даёт~\eqref{eq:bound}.
Собственный вектор с собственным значением нуль обеспечивает остроту.
Формула~\eqref{eq:residual} разворачивает эту оценку по координатам;
условный перенос ниже получается подстановкой заданных гипотез.
Это специализация известного спектрального принципа, а не новая
связь между полными алгебрами Альберта и Грисса.

В частности,
\begin{align}
 S(R_{1/4}x)&=\tfrac1{16}S(x)+\tfrac12x_0^2,\label{eq:quarter}\\
 S(R_{1/32}x)&=\tfrac1{1024}S(x)+\tfrac{15}{16}x_0^2
 +\tfrac7{16}\bigl(|u_{01}|_2^2+|u_{02}|_2^2\bigr).
 \label{eq:thirtysecond}
\end{align}
Равенство в соответствующей нижней оценке при $1/4$ эквивалентно
$x_0=0$. При $1/32$ дополнительно должна обращаться в нуль сумма
$|u_{01}|_2^2+|u_{02}|_2^2$. Эти условия также проверены в Lean.

\begin{corollary}[Условный перенос оценки]\label{cor:map}
Пусть множество $A$ снабжено умножением и действием рациональных
скаляров. Пусть $e,x\in A$, $f\colon A\to\J$ и
$\lambda\in\QQ$ удовлетворяют
\begin{equation}\label{eq:maphyp}
 ex=\lambda x,\qquad f(e)=P,\qquad f(\lambda x)=\lambda f(x).
\end{equation}
Если $\lambda\le1/4$, то
\[
 S\bigl(f(e)\circ f(x)-f(ex)\bigr)\ge\lambda^2S(f(x)).
\]
Если дополнительно $\lambda\ne0$ и $S(f(x))>0$, то
$f(e)\circ f(x)\ne f(ex)$.
\end{corollary}
\begin{proof}
По~\eqref{eq:maphyp} разность внутри $S$ равна $R_\lambda f(x)$.
Применим теорему~\ref{thm:residual}. При дополнительных условиях
правая часть положительна, тогда как $S(0)=0$.
\end{proof}
\leanref{SpectralDefect.lean}{productError_lower_bound; product_not_preserved_of_positive_size}

Глобальная линейность $f$ и аксиомы алгебры для $A$ в этой формулировке
не требуются: достаточно явно указанных данных. Числа $1/4$ и $1/32$
мотивированы известным спектром Monster-типа~\cite[Example~2.5]{mcinroy}.
Однако полная алгебра Грисса, Monster VOA и конкретная карта из них
в $\J$ здесь не построены. Острота теоремы~\ref{thm:residual} относится
к остаткам в целевом пространстве, а не к классу таких отсутствующих карт.

\subsection{Конкретный пример и граница согласования}
Для $Y=\diag(0,1,0)$ и $\lambda=1/4$ непосредственно получаем
\[
 P\circ Y=0,\qquad R_{1/4}Y=-\tfrac14Y,\qquad
 S(Y)=1,\qquad S(R_{1/4}Y)=\tfrac1{16}.
\]
Это тот же свидетель остроты, который проверен в Lean.
Если в исходной модели $ex=\tfrac14x$, а отображение удовлетворяет
$f(e)=P$, $f(x)=Y$ и $f(\tfrac14x)=\tfrac14Y$, то
\[
 f(e)\circ f(x)-f(ex)=-\tfrac14Y\ne0.
\]
Следовательно, такое отображение не может сохранять это произведение.
Пример задаёт проверяемое условие на отображение и не постулирует
построение конкретной алгебры-источника.

\subsection{Какая общая структура выделена}
Общими данными второй пары являются три ортогональных блока
размерностей $1$, $16$, $10$ с весами действия $1$, $1/2$, $0$
и заданными квадратичными вкладами. В объекте $\mathcal A_2$ они
возникают из умножения, а в $\mathcal B_2$ заданы спектрально.
Соответствие $U$ согласует эти два представления.

Эта взвешенная структура пространства является кандидатом на
суперобъект второй пары. Доказанная формула остатка показывает,
какую количественную информацию сохраняет соответствие.
Применение к внешней модели даёт лишь условный критерий
несовместимости отображения при гипотезах следствия~\ref{cor:map};
конкретное объединяющее вложение внешней алгебры не построено.

\section{Третья пара: отношение на тройках и геометрический код}\label{sec:code}
\subsection{Определения двух объектов}
Стандартный мнимый базис в конвенции~\eqref{eq:oct} обозначим
\[
 v_1=(i,0),\ v_2=(j,0),\ v_3=(k,0),\ v_4=(0,1),
 \ v_5=(0,i),\ v_6=(0,j),\ v_7=(0,k).
\]
Через $T$ обозначим множество 35 неупорядоченных троек различных
меток из $\{1,\ldots,7\}$.

\begin{definition}[Отношение нулевых ассоциаторов]\label{def:incidence}
Первый объект --- семейство отношений $\mathcal A_3=\{I_p:p\in S_7\}$
на $T$. Для $i<j<k$ тройка $\{i,j,k\}$ принадлежит $I_p$, если
\[
 (v_{p(i)}v_{p(j)})v_{p(k)}-v_{p(i)}(v_{p(j)}v_{p(k)})=0.
\]
Перестановка $p$ меняет базисные метки; совпадающие отношения
считаются одним элементом семейства.
\end{definition}

\begin{definition}[Код размеченных плоскостей Фано]\label{def:fano-code}
Пусть $F$ --- плоскость Фано на метках $1,\ldots,7$ с линиями
$123,145,167,246,257,347,356$. Для $p\in S_7$ положим
\[
 F_p=\{t\in T:p(t)\in F\},\qquad
 b_p(t)=\begin{cases}1,&t\in F_p,\\0,&t\notin F_p.\end{cases}
\]
Второй объект --- двоичный код $\mathcal B_3=\{b_p:p\in S_7\}$
длины 35, образованный всеми такими перемаркировками плоскости.
Здесь $p(t)$ означает применение $p$ ко всем меткам тройки.
\end{definition}

\subsection{Какую связь устанавливаем}
Отношению $I_p$ сопоставим его характеристическое слово $c_p$:
\begin{equation}\label{eq:incidence}
 c_p(\{i,j,k\})=1\quad\Longleftrightarrow\quad
 (v_{p(i)}v_{p(j)})v_{p(k)}-v_{p(i)}(v_{p(j)}v_{p(k)})=0,
 \quad i<j<k.
\end{equation}
Требуется установить $c_p=b_p$ для каждой перестановки.
Тогда отношение нулевых ассоциаторов и геометрическое кодовое слово
описывают одну и ту же инцидентность. Далее проверяется более сильное
свойство этой связи: полное отношение восстанавливается по части
координат слова, в том числе при одном ошибочном ответе.

\subsection{Почему два определения задают одну инцидентность}
В используемой конвенции для различных меток ассоциатор равен нулю
тогда и только тогда, когда побитовое исключающее ИЛИ их числовых
значений равно нулю:
\[
 i\mathbin{\oplus}j\mathbin{\oplus}k=0.
\]
Полный критерий --- совпадение хотя бы двух меток или нулевое
исключающее ИЛИ --- проверен на конкретном произведении~\eqref{eq:oct}
для всех $7^3=343$ упорядоченных троек, включая повторяющиеся метки.
Отсюда для различных меток получается известная инцидентность
плоскости Фано~\cite{baez}. Для выбранных меток её линии совпадают
с $F$, поэтому $I_p=F_p$ и $c_p=b_p$ для каждой $p$.

Множество $C=\{c_p:p\in S_7\}$ содержит 30 различных слов длины 35.
В каждом слове семь единиц, в каждом столбце каталога --- шесть.
Здесь $C$ --- конечное нелинейное множество; оно не отождествляется
с линейным $(35,14)$-кодом, изученным в~\cite{liebler}.
В формализации каталог предъявлен явно и проверен против вычисленных
представителей. Отдельно проверено, что он покрывает все 5040
перестановок. Тем самым результаты относятся ко всему $C$,
а не только к произвольно выбранным примерам.

Для упорядоченного набора тестов $Q=(t_1,\ldots,t_m)$ обозначим
через $c|_Q=(c(t_1),\ldots,c(t_m))$ его ограничение.
Расстояние $\Ham$ есть число несовпадающих битов.
Таким образом, общий объект --- семейство $C$ с отображениями
наблюдения $c\mapsto c|_Q$. Геометрическая линия соответствует
нулевому ассоциатору и единичному биту; ограничение слова соответствует
выбору алгебраических тестов. Следующая теорема доказывает, что
часть этих наблюдений достаточна для восстановления всей структуры.

\subsection{Утверждение и доказательство}

\begin{theorem}[Сокращённые сертификаты]\label{thm:probes}
Определим наборы тестов, записывая тройку $\{i,j,k\}$ как $ijk$:
\begin{align*}
 Q_{10}={}&(123,127,135,137,157,236,256,267,356,567),\\
 Q_{18}={}&(124,126,145,147,156,167,234,236,245,\\
          &\phantom{(}247,256,267,345,347,356,367,457,567).
\end{align*}
Тогда выполняются следующие утверждения.
\begin{enumerate}
 \item Отображение $c\mapsto c|_{Q_{10}}$ инъективно на $C$.
 \item Для любых различных $c,c'\in C$ имеем
 $\Ham(c|_{Q_{18}},c'|_{Q_{18}})\ge3$.
 Следовательно, если полученное слово $r\in\{0,1\}^{18}$ находится
 на расстоянии не более одного от $c|_{Q_{18}}$ для некоторого
 $c\in C$, то такое $c$ единственно.
 \item Ни один заранее фиксированный набор из не более чем восьми
 тестов не задаёт инъективного ограничения на $C$.
\end{enumerate}
Таким образом, минимальное число фиксированных безошибочных тестов
$N_{\min}$ удовлетворяет $9\le N_{\min}\le10$.
\end{theorem}
\begin{proof}
Первые два утверждения о сокращённых словах проверяются конечным
сертификатом: для всех 30 элементов проверенного каталога
вычислены ограничения на указанные позиции (таблица~\ref{tab:words}); проверены инъективность
первого ограничения и нижняя граница расстояния для всех различных
пар второго. Эти проверки выполнены ядром Lean в точной арифметике.
Если бы два различных слова находились на расстоянии не более одного
от $r$, неравенство треугольника дало бы расстояние между их
ограничениями не больше двух, вопреки проверенной границе три.

Для третьего утверждения предположим, что восемь фиксированных
столбцов различают все 30 слов. Их суммарный вес по всем строкам
равен $8\cdot6=48$. С другой стороны, среди различных восьмибитных
слов существует только одно слово веса нуль и восемь слов веса один.
Поэтому суммарный вес любых 30 различных слов не меньше
\[
 0+8\cdot1+21\cdot2=50,
\]
что противоречит предыдущему равенству. Если число тестов положительно
и меньше восьми, дополним их до восьми повторениями: инъективность
сохранилась бы, а сумма каждого столбца по-прежнему равнялась бы шести.
При нуле тестов все ограничения совпадают. Следовательно, запрещены
все числа тестов от нуля до восьми. Первый пункт даёт верхнюю границу 10.
\end{proof}
\leanref{OctonionIncidence.lean}{ten_probes_reconstruct; eighteen_probes_reconstruct_with_one_error}
\leanref{OctonionIncidence.lean}{at_most_eight_fixed_probes_insufficient}

\paragraph{Сопоставление с литературой.}
Нижняя граница $N_{\min}\ge9$ является частным случаем
\cite[Proposition~3]{crowston}: если имеется $n$ объектов и каждый
двоичный тест положителен не более чем на $r$ объектах, то для
различения всех объектов требуется не менее
$\lceil2(n-1)/(r+1)\rceil$ фиксированных тестов.
Здесь объекты --- 30 слов $C$, а каждый запрос положителен на шести
словах; следовательно,
\[
 N_{\min}\ge\left\lceil\frac{2(30-1)}{6+1}\right\rceil=9.
\]
Это применение известной границы, а не новая общая нижняя оценка.
Наборы $Q_{10}$ и $Q_{18}$ проверены как конкретные сертификаты.
В просмотренных первичных источниках не установлен их точный
или изоморфный предшественник; историческая новизна этих наборов
остаётся неизвестной. Наличие таких сертификатов само по себе
не устанавливает их оптимальности или новизны связи
<<октонионы--плоскость Фано--код>>.

\subsection{Конкретный пример и смысл связи}
Для исходного, непереставленного базиса семь положительных троек имеют вид
\[
 123,\quad145,\quad167,\quad246,\quad257,\quad347,\quad356.
\]
Они задают первую строку каталога в таблице~\ref{tab:words}.
Одни и те же данные прочитываются как семь нулевых ассоциаторов,
семь линий размеченной плоскости Фано и семь единиц слова длины 35.
В порядке тестов теоремы~\ref{thm:probes} получаем
\[
 c|_{Q_{10}}=\texttt{1000000010},\qquad
 c|_{Q_{18}}=\texttt{001001000000011000}.
\]
Первое слово однозначно определяет все семь троек среди 30 вариантов.
Во втором слове изменим ответ на первый тест $124$ с нуля на единицу:
\[
 r=\texttt{101001000000011000}.
\]
Исходное слово находится от $r$ на расстоянии один. Другого
кодового слова в этом радиусе нет по теореме~\ref{thm:probes}.
Следовательно, даже ошибочный набор наблюдений однозначно определяет
исходную инцидентность при условии не более одной ошибки.

\subsection{Какая общая структура выделена}
Доказанное совпадение $I_p=F_p$ выделяет общее отношение инцидентности:
в первом объекте оно задано нулевым ассоциатором, во втором ---
геометрической линией и единичным битом. Семейство таких отношений
вместе с отображениями наблюдения $c\mapsto c|_Q$ является
кандидатом на суперобъект третьей пары.

Теорема~\ref{thm:probes} показывает содержательность этой структуры:
часть наблюдений определяет всё отношение, а достаточное расстояние
между словами обеспечивает единственность при одной ошибке.
Таким образом, из связи двух определённых объектов получены
общие данные и проверенные правила их восстановления.

\subsection{Как устроена конечная часть доказательства}
Каталог и таблицы не принимаются как внешние аксиомы.
Перечисление всех перестановок разбито на 50 блоков по 100 элементов
и последний блок из 40 элементов. Для каждого блока проверяется
принадлежность вычисленных слов каталогу. Общая лемма о разбиении
списка объединяет 51 сертификат в покрытие всего перечисления.
Используется \texttt{decide +kernel}; независимый Python-эксперимент
не входит в логическую зависимость этих доказательств.

Восстановление в теореме~\ref{thm:probes} означает восстановление
35-битной инцидентности. Оно не восстанавливает ориентацию,
знаки таблицы умножения, произвольный непрерывно изменённый базис
или индекс перестановки среди 5040 вариантов. Для 18 тестов
доказана единственность при обещании не более одной ошибки;
отдельная реализация алгоритма-декодера в Lean не построена.
Вопрос о девяти тестах и оптимальность числа 18 в настоящей работе
не решены.

\section{Формальная проверка и воспроизводимость}\label{sec:verification}
Проверка выполнена в Lean~4.19.0 с mathlib, зафиксированной на коммите
\nolinkurl{c44e0c8ee63ca166450922a373c7409c5d26b00b}.
Полная проверка охватывает 11 исходных Lean-файлов и
149 деклараций теорем, включая вспомогательные леммы
и 51 блок конечного сертификата. Результат запуска для
публикуемой версии, коды завершения и контрольные суммы
приведены в \nolinkurl{checks/last-run-summary.json}.

Обёртка \nolinkurl{scripts/check.py} запускает проверку в чистой копии;
сценарий \nolinkurl{scripts/verify.py} сопоставляет рекурсивный список
Lean-исходников со списком запуска. Для каждого файла
сохранены команда, код завершения, SHA-256 исходника и результат
сравнения хешей до и после компиляции, хеш результата \texttt{.olean}
и хеш журнала.
Импортированные кеши mathlib использованы повторно; вся библиотека
заново не компилировалась.

Для каждой декларации выполнена команда \texttt{\#print axioms}.
Архитектурный контекст проверяющих термов Lean описан в~\cite{demoura};
конкретная версия и успешность этого запуска удостоверяются журналами проверки.
Среди разрешённых оснований встречаются только
\texttt{propext}, \texttt{Classical.choice} и \texttt{Quot.sound}.
Пропуски доказательств \texttt{sorry}, пользовательские аксиомы
и \texttt{native\_decide} в проверенных исходниках отсутствуют.
Доверительная база включает используемое ядро Lean, импортированные
определения и указанные основания; отдельная реализация проверяющего
ядра не запускалась.

\subsection{Связь текста и формального артефакта}
Математические доказательства выше являются изложением формальных
результатов, а не автоматическим переводом доказательных термов.
Ядро проверяет точные определения и утверждения в Lean-файлах;
оно не проверяет этот естественно-языковой текст, историческую
новизну или соответствие общей исследовательской цели.
Смысловая сверка формулировок выполнялась программными ИИ-агентами.
Подтверждённой внешней рецензии математиком-человеком нет.

Проверка построена как последовательность: точная постановка,
сопоставление с литературой, вычислительный эксперимент,
формальное доказательство и сверка его смысла с текстом.
Численные примеры выше специализируют общие теоремы; они не
заменяют доказательства для произвольных допустимых параметров.
Само название <<суперобъект>> не является отдельным предикатом Lean.
Формально подтверждаются конкретные тождества, восстановления
и препятствия, которыми определена связь представлений.
Шесть определений в тексте организуют эти результаты в три пары.
Отображение $U$ во второй паре является явной записью перегруппировки
координат: его свойства выведены из проверенных формул действия
$M_P$ и формы $S$. Отдельная Lean-декларация для $U$ и абстрактное
определение суперобъекта в формальный артефакт не добавлялись.

\section{Выделенные общие структуры и границы вывода}
\paragraph{Первая пара.}
Семейство операторов и геометрия размеченных троек связаны через
веса расстояний. Доказанное восстановление устанавливает, что эти
веса определяют оба представления. Кандидат на суперобъект ---
размеченная структура весов с операторным, геометрическим
и графовым представлениями.

\paragraph{Вторая пара.}
Пространство с оператором умножения и трёхблочная спектральная модель
связаны перегруппированием координат, сохраняющим выделенный оператор
и квадратичную форму. Кандидат на суперобъект --- взвешенное
ортогональное разложение пространства. Формула энергии остатка
и её острая оценка описывают сохраняемую количественную информацию.

\paragraph{Третья пара.}
Отношение нулевых ассоциаторов и геометрический код связаны
характеристическим словом одной инцидентности. Кандидат на суперобъект
--- семейство инцидентностей с отображениями наблюдения.
Доказано восстановление по 10 тестам и единственность восстановления
по 18 тестам при не более чем одной ошибке.

Во всех трёх парах можно явно выделить общие данные и указать,
какие представления они согласуют. Именно это обосновывает
выделение кандидатов на суперобъекты в принятом рабочем смысле.
Доказательство высокой внутренней сложности или существования
единого объекта, объединяющего все три пары, потребовало бы
дополнительной постановки и отдельных результатов.

Дальнейшие математические вопросы --- существование девяти
безошибочных тестов, минимальность числа устойчивых тестов и
построение конкретных исходных моделей для переноса произведения.
Статус исторической новизны сохраняется в объёме литературного
сопоставления: известные основания отделены от конкретных наборов
тестов, приоритет которых не установлен.

\appendix
\section{Состав формального приложения}\label{app:files}
В таблице приведены все файлы последнего полного запуска.
Число деклараций не равно числу самостоятельных исследовательских
результатов. Имена модулей указаны без общего префикса; точные пути
даны в сертификате проверки. Общий модуль импорта не имеет
собственных теорем.
\begin{center}\small
\begin{tabular}{@{}p{.73\textwidth}rr@{}}
\toprule
Файл & Декларации & Код\\
\midrule
\nolinkurl{MinimalAInfinity.lean} & 2 & 0\\
\nolinkurl{AlbertObstruction.lean} & 4 & 0\\
\nolinkurl{AlbertMatrix.lean} & 11 & 0\\
\nolinkurl{Derivations.lean} & 7 & 0\\
\nolinkurl{AssociativeBridge.lean} & 4 & 0\\
\nolinkurl{TripleBridge.lean} & 2 & 0\\
\nolinkurl{ProbeLowerBound.lean} & 5 & 0\\
\nolinkurl{AlbertSpectralGeometry.lean} & 18 & 0\\
\nolinkurl{SpectralDefect.lean} & 27 & 0\\
\nolinkurl{OctonionIncidence.lean} & 69 & 0\\
Общий модуль импорта & 0 & 0\\
\midrule
Всего & 149 & \\
\bottomrule
\end{tabular}
\end{center}

Помимо трёх основных постановок, приложение содержит общие леммы
о дифференцированиях, универсальной обёртывающей алгебре, тройных
системах и препятствии минимальному бинарному произведению.
Эти леммы не подменяют отсутствующие полные конструкции $E_{10}$,
Monster VOA или градуированной $A_\infty$-теории.

\section{Сокращённые слова конечного сертификата}
Нумерация строк соответствует \nolinkurl{catalogueWords} в Lean-исходнике.
Порядок битов совпадает с порядком тестов $Q_{10}$ и $Q_{18}$
в теореме~\ref{thm:probes}. Все строки первого столбца слов различны;
у второго столбца слов попарные расстояния не меньше трёх.
Таблица является извлечением проверенного каталога, а доказательство
его полноты относительно $S_7$ хранится в формальном артефакте.
\begingroup\small\renewcommand{\arraystretch}{1.02}
\begin{longtable}{@{}rll@{}}
\caption{Ограничения 30 слов инцидентности.}\label{tab:words}\\
\toprule
№ & Слово из 10 битов & Слово из 18 битов\\
\midrule\endfirsthead
\toprule
№ & Слово из 10 битов & Слово из 18 битов\\
\midrule\endhead
\bottomrule\endfoot
1 & \texttt{1000000010} & \texttt{001001000000011000} \\
2 & \texttt{1000001000} & \texttt{001001000110000000} \\
3 & \texttt{1000100110} & \texttt{000000001001011000} \\
4 & \texttt{1000101000} & \texttt{000000000110100100} \\
5 & \texttt{1000000100} & \texttt{000110001001000000} \\
6 & \texttt{1000000000} & \texttt{000110000000100100} \\
7 & \texttt{0010010000} & \texttt{100001010000010000} \\
8 & \texttt{0010001000} & \texttt{100001000010000010} \\
9 & \texttt{0000100100} & \texttt{100000000001010000} \\
10 & \texttt{0000101000} & \texttt{100000000010100000} \\
11 & \texttt{0001000100} & \texttt{100010000001000010} \\
12 & \texttt{0001010000} & \texttt{100010010000100000} \\
13 & \texttt{0000010000} & \texttt{000001010100000000} \\
14 & \texttt{0000000010} & \texttt{000001000000001010} \\
15 & \texttt{0000000100} & \texttt{000100100001000000} \\
16 & \texttt{0000000001} & \texttt{000100000000100001} \\
17 & \texttt{0001000110} & \texttt{000000100001001010} \\
18 & \texttt{0001010001} & \texttt{000000010100100001} \\
19 & \texttt{0000100000} & \texttt{010000000100000100} \\
20 & \texttt{0000100010} & \texttt{010000001000001000} \\
21 & \texttt{0010000000} & \texttt{010100100000000100} \\
22 & \texttt{0010000001} & \texttt{010100001000000001} \\
23 & \texttt{0001000010} & \texttt{011000100000001000} \\
24 & \texttt{0001000001} & \texttt{011000000100000001} \\
25 & \texttt{0100000000} & \texttt{000010000000000110} \\
26 & \texttt{0100010000} & \texttt{000010011000000000} \\
27 & \texttt{0110001000} & \texttt{000000100010000110} \\
28 & \texttt{0110010001} & \texttt{000000011000010001} \\
29 & \texttt{0100001000} & \texttt{001000100010000000} \\
30 & \texttt{0100000001} & \texttt{001000000000010001} \\
\end{longtable}
\endgroup

\section{Артефакты проверки}
Публичный комплект расположен в репозитории
\url{https://github.com/Medrov/astra-math-research}.
Lean-модули находятся в каталоге \nolinkurl{MathResearch/},
общий модуль импорта --- \nolinkurl{MathResearch.lean}.
Формальные определения и утверждения этой редакции сохраняют
математическое содержание предшествующего препринта.

Результат контрольного запуска для опубликованных исходников
сохранён в \nolinkurl{checks/last-run-summary.json}.
Он содержит версию Lean, ревизию mathlib, состав проверки,
коды завершения, зависимости аксиом и контрольные суммы исходников.
Журналы отдельных модулей находятся в
\nolinkurl{checks/lean-logs/}. Для повторной проверки выполните
\mbox{\texttt{python3\ scripts/check.py}} из корня репозитория;
новый локальный результат появится в
\nolinkurl{verification/verification.json}.

Порядок подготовки окружения описан в \nolinkurl{REPRODUCE.md},
а работа агентов на этих трёх задачах --- в
\nolinkurl{WALKTHROUGH.md} и \nolinkurl{docs/Pipeline.pdf}.
Интерактивные примеры приведены в \nolinkurl{docs/Interactive.html}.
Литературное сопоставление и его границы изложены в основном тексте
и библиографии; повторная компиляция не устанавливает историческую
новизну результатов.

\paragraph{Статус рукописи.}
Рабочий препринт; не рецензирован и не размещён
в научном архиве. Подготовка текста и формальных артефактов
выполнялась с использованием ИИ-агентов. Персональный авторский
состав и аффилиации в этой версии не указаны.

\begin{thebibliography}{99}
\bibitem{gowda} M.~Seetharama Gowda. \emph{Some majorization inequalities induced by Schur products in Euclidean Jordan algebras}. arXiv:1911.00579v1, 2019. \url{https://arxiv.org/abs/1911.00579v1}.
\bibitem{gowda2012} M.~Seetharama Gowda, J.~Tao, R.~Sznajder. Complementarity properties of Peirce-diagonalizable linear transformations on Euclidean Jordan algebras. \emph{Optimization Methods and Software}, 27(4--5):719--733, 2012; online 2011. \url{https://doi.org/10.1080/10556788.2011.626036}.
\bibitem{dokmanic} I.~Dokmani\'c, R.~Parhizkar, J.~Ranieri, M.~Vetterli. Euclidean Distance Matrices: Essential Theory, Algorithms and Applications. \emph{IEEE Signal Processing Magazine}, 32(6):12--30, 2015. \url{https://doi.org/10.1109/MSP.2015.2398954}. arXiv:1502.07541v2.
\bibitem{kalofolias} V.~Kalofolias. How to Learn a Graph from Smooth Signals. \emph{Proceedings of Machine Learning Research}, 51:920--929, 2016. \url{https://proceedings.mlr.press/v51/kalofolias16.html}.
\bibitem{pankrashkin} K.~Pankrashkin. \emph{Introduction to spectral theory}. Lecture notes, 2014, Theorem~6.1, p.~51. \url{https://www.imo.universite-paris-saclay.fr/~nonnenmacher/enseign/Spectral_Theory_notes_Pankrashkin_2014.pdf}.
\bibitem{mcinroy} J.~McInroy, S.~Shpectorov. \emph{Axial algebras of Jordan and Monster type}. arXiv:2209.08043v1, 2022. \url{https://arxiv.org/abs/2209.08043v1}.
\bibitem{baez} J.~C.~Baez. The Octonions. \emph{Bulletin of the American Mathematical Society}, 39(2):145--205, 2002. arXiv:math/0105155v4. \url{https://arxiv.org/abs/math/0105155v4}.
\bibitem{liebler} R.~A.~Liebler, K.-H.~Zimmermann. Combinatorial $S_n$-Modules as Codes. \emph{Journal of Algebraic Combinatorics}, 4(1):47--68, 1995. \url{https://doi.org/10.1023/A:1022485624417}.
\bibitem{crowston} R.~Crowston, G.~Gutin, M.~Jones, G.~Muciaccia, A.~Yeo. \emph{Parameterizations of Test Cover with Bounded Test Sizes}. arXiv:1209.6528v2, 2013 (первая версия 2012). \url{https://arxiv.org/abs/1209.6528v2}.
\bibitem{demoura} L.~de Moura, S.~Kong, J.~Avigad, F.~van Doorn, J.~von Raumer. The Lean Theorem Prover (System Description). In: \emph{Automated Deduction---CADE-25}, LNCS 9195, pp.~378--388. Springer, 2015. \url{https://doi.org/10.1007/978-3-319-21401-6_26}.
\end{thebibliography}

\end{document}
